% Copyright (c) 2026 Benjamin Stanley Frohman. All rights reserved.
% Released under Apache 2.0 license as described in the file LICENSE.
% Authors: Benjamin Stanley Frohman (ORCID: 0009-0006-7068-3718)
% Origin Provenance: April 2026 (Compiled 2026-07-04 15:59:24 CDT; Google Drive Archive 2026-07-07 14:20:55 CDT)
\documentclass[11pt,a4paper]{article}
\usepackage[utf8]{inputenc}
\usepackage[T1]{fontenc}
\usepackage{amsmath,amssymb,amsthm}
\usepackage{geometry}
\geometry{margin=1in}

\title{A Note on the Frohmanian Symplectic Tether Framework Applied to a Functional Equation}
\author{Benjamin Frohman}
\date{Origin Provenance: April 2026 (Compiled July 4, 2026; Archived July 7, 2026) \\ \small ORCID: \texttt{0009-0006-7068-3718} \quad|\quad Licensed under Apache-2.0}

\begin{document}

\maketitle

\begin{abstract}
We illustrate how the methodological tools of the Frohmanian Symplectic Tether Theorem --- uniqueness via canonicity conditions, invariance under iteration (coadjoint analogue), reduced-orbit necessity, and independent majorant control --- directly yield clean first-principles proofs of the three statements in problem B4 concerning continuous real-valued functions \(f\) satisfying \(f(x^2) = f(x)^2\).
\end{abstract}

\section{Setup}
Let \(f: D \to \mathbb{R}\) be continuous and satisfy the multiplicative invariance
\[
f(x^2) = f(x)^2 \quad \forall x \in D.
\]
Define the (forward/backward) orbit of a point \(x\) under the iteration map \(\phi(y) = y^2\). By continuity, behavior along orbits is controlled. This invariance is analogous to coadjoint invariance in the \(\mathfrak{sdiff}(\mathbb{T}^3)\) setting.

\section{Proof of (i): Uniqueness on \([0,1]\) with \(f(0) \neq 0\)}

Consider any \(x \in [0,1]\). The forward orbit satisfies \(x^{2^n} \to 0\). Then
\[
f(x^{2^n}) = f(x)^{2^n}.
\]
Taking \(n \to \infty\), continuity implies
\[
f(0) = \lim_{n\to\infty} f(x)^{2^n}.
\]
Assume \(f(0) \neq 0\). Suppose there exists \(x\) with \(f(x) \neq 1\):
- If \(|f(x)| > 1\), then \(|f(x)^{2^n}| \to \infty\), contradicting continuity at 0.
- If \(|f(x)| < 1\), then \(f(x)^{2^n} \to 0\), again contradicting \(f(0) \neq 0\).

Thus \(f(x) = 1\) (sign consistency follows from squaring). By density of orbits in \([0,1]\) under iteration and continuity, \(f \equiv 1\) on \([0,1]\).

This mirrors the **5-step canonicity** (C1--C5) and uniqueness theorem: the invariance forces a unique canonical solution, with degeneracy on constant proxies ruling out alternatives.

\section{Proof of (ii): All-or-Nothing Non-Vanishing on \((0,\infty)\)}

Suppose there exists \(a > 0\) with \(f(a) = 0\). Then \(f(a^{1/2^n}) = 0\) for all \(n\). The backward orbit \(a^{1/2^n} \to 1\), so continuity yields \(f(1) = 0\).

From \(f(1) = 0\), forward iteration and density of orbits under roots/powers fill \((0,\infty)\) by continuity. Thus \(f \equiv 0\).

Conversely, non-vanishing at one point propagates everywhere.

This is the direct analogue of **reduced-orbit necessity** (T3 shear counterexample): without the proper reduction (connected components under the iteration map), global invariance fails.

\section{Proof of (iii): Infinitely Many \(f\) on \((0,\infty)\) with \(\int_0^\infty f(x)\,dx < 1\)}

On \((1,\infty)\), prescribe \(f\) freely on a fundamental domain (intervals not in the image of squaring from below) and extend via \(f(x) = \pm \sqrt{f(x^2)}\) (continuous branch). Choose decay near 0 and \(\infty\) (e.g., \(f(x) \sim x^r\) for \(r>0\), or localized bumps) so the integral is finite and arbitrarily small.

The freedom on the fundamental domain yields infinitely many distinct solutions.

This follows the **independent majorant + negative quadratic feedback** philosophy: the squaring relation plus continuity provides controllable decay (quartic absorption analogue), while the tether proxy in the discrete setting mirrors the freedom under global constraints.

\section{Conclusion: Direct Methodological Transfer}
The Frohmanian framework supplies the precise toolkit:
\begin{itemize}
\item Uniqueness/canonicity under forcing conditions (parts (i)--(ii)).
\item Orbit reduction and invariance (coadjoint analogue).
\item Independent majorant control for global bounds (part (iii)).
\end{itemize}
No ansatz is required; all results follow from minimal geometric/functional constraints + continuity, exactly as in the tether derivation for global Navier--Stokes regularity.

\end{document}